\section{Proofs and exact certificates for depth bounds}\label{dp:appendix}

\subsection{Orbit equations and vertical completion}

For $G=\left(\begin{smallmatrix}g_{11}&g_{12}\\g_{21}&g_{22}\end{smallmatrix}\right)$
with positive determinant, a direct $2\times2$ identity gives
\begin{equation}\label{dp:orbit-equation}
 C_G=\left\{s:\det G\,q(s)\ge
 \frac{(g_{11}x-g_{22}y-g_{21}w+g_{12})^2}{4},\quad
 g_{11}w+g_{12}y+g_{21}x+g_{22}\ge0\right\}.
\end{equation}
Indeed $\det\sym A=\det A-(A_{12}-A_{21})^2/4$, and a symmetric
$2\times2$ matrix is PSD if and only if its determinant and trace are
nonnegative. This describes a trace-selected branch of the scalar quadratic
region. It is not a classification by connected components: two branches
can meet at an apex.

\begin{lemma}[Exact completion membership]\label{dp:completion-membership}
The sum $B_G=C_G+\R_+e_w$ is closed. A point $s$ belongs to it if and only if
$q(s)\ge0$ and
\begin{equation}\label{dp:lowering}
 \sym\bigl(GM(s-\tau e_w)\bigr)\succeq0
 \quad\text{for some }0\le\tau\le q(s).
\end{equation}
\end{lemma}

\begin{proof}
\Cref{dp:orbit-equation} implies $q\ge0$ on $C_G$.
If $s=c+\tau e_w$, then $q(c)=q(s)-\tau\ge0$, giving the stated
bound on $\tau$. Conversely a lowered point proves membership.
If $s_n=c_n+\tau_ne_w\to s$, then $0\le\tau_n\le q(s_n)$ is bounded.
A subsequence converges to $\tau\ge0$, and the closed PSD inequality gives
$s-\tau e_w\in C_G$. Thus the sum is closed.
\end{proof}

\subsection{Analytic collapse of the completed family}\label{dp:vanishing-proof}

We prove the final assertion of \cref{dp:vanishing}. Suppose it fails.
Then, after reducing a positive lower bound if necessary, there are
$0<z<1$, $\varepsilon_n\downarrow0$, and admissible completed sets
$B_{G_n}$ containing
\[
 T_n=\conv\{\sbar_n,\sbar_n+zp_1,\sbar_n+zp_2,\sbar_n+zp_3\},
 \qquad\sbar_n=(0,0,\varepsilon_n).
\]
Normalize $\norm{G_n}_F=1$ and take a subsequence with $G_n\to G\ne0$.
The determinant of $G$ is nonnegative. By
\cref{dp:completion-membership}, each vertex has a lowered point in
$C_{G_n}$ with lowering between zero and its $q$ value. Pass to limits of
these four bounded lowerings. The limiting lowered points are
\[
 \ell_0=0,\quad\ell_1=(z,-z,-\tau_1),\quad
 \ell_2=(z,-2z,-\tau_2),\quad\ell_3=(z,z,z-\tau_3),
\]
where $0\le\tau_1\le z^2$, $0\le\tau_2\le2z^2$, and
$0\le\tau_3\le z-z^2$.

First suppose $\det G>0$. The limiting completed set $B_G$ contains
the full-dimensional tetrahedron
$T_0=\conv\{0,zp_1,zp_2,zp_3\}$. It is $S$-free and contains zero.
It must lie in $w\ge0$. Otherwise a point $b\in B_G$ with $b_w<0$
gives $q(\theta b)<0$ for sufficiently small positive $\theta$.
Taking a convex combination of $\theta b$ with an interior point of
$T_0$, with weight sufficiently close to one on $\theta b$, produces
an interior point of $B_G$ in $\{q<0\}$, a contradiction.
Consequently $C_G\subseteq B_G$ also lies in $w\ge0$. A point of
$B_G$ with $w=0$ cannot have been lowered by a positive amount, and
therefore belongs to $C_G$.

In particular $0,zp_1,zp_2\in C_G$. At zero, PSD membership forces
the second column of $G$ to be a positive multiple of $e_2$.
After positive scaling,
$G=\left(\begin{smallmatrix}\alpha&0\\\beta&1\end{smallmatrix}\right)$
with $\alpha>0$. At $w=0$, the first diagonal entry of
$\sym(GM)$ is zero, so PSD membership forces its off-diagonal entry
$(\alpha x+y)/2$ to vanish. The two horizontal vertices require both
$\alpha=1$ and $\alpha=2$, a contradiction.

Now suppose $\det G=0$. Write $G=a b^T$ with $a,b\ne0$.
The PSD condition for a rank-one product gives
$M(\ell_i)^Tb\in\R_+a$. In particular all $\ell_i$ lie in a
nontrivial affine plane. The equation is nontrivial because the condition
that $M(s)^Tb$ be parallel to a fixed nonzero $a$ cannot hold identically
in the three coordinates unless $b=0$.
Their planar projections are not collinear, so this plane has the form
$w=Ax+By$. Let
\[
 U=\conv\{\ell_i:0\le i\le3\}+\R_+e_w.
\]
This full-dimensional set is $S$-free. To see the limiting step explicitly,
the corresponding lowered points before taking limits generate subsets
$U_n\subseteq B_{G_n}$. A point of $\intt U$ has an interior planar
projection and positive slack above its lower plane. Expressing its
projection with strictly positive convex weights on the limiting vertices,
then perturbing those weights and using its positive vertical slack, shows
that a neighborhood of the point belongs to $U_n$ for all sufficiently large
$n$. If that point also had $q<0$, $B_{G_n}$ would fail to be $S$-free.
Thus $Ax+By\ge xy$ on the projection
\[
 \Delta=\conv\{(0,0),(z,-z),(z,-2z),(z,z)\}.
\]
Interior projections first give this inequality, and continuity gives it
on the boundary. Near zero along $(1,-1)$ and $(1,-2)$ it yields
$A-B\ge0$ and $A-2B\ge0$. The two lowered horizontal vertices give
$A-B=-\tau_1/z\le0$ and $A-2B=-\tau_2/z\le0$. Hence $A=B=0$.
Then $\ell_3$ lying in the plane requires $\tau_3=z$, contrary to
$\tau_3\le z-z^2$. Both cases are impossible. This proves vanishing.

\subsection{The finite rational certificate for the upper constants}\label{dp:box-proof}

The upper bounds in \cref{dp:certified-sharpness} have one common
certificate scheme. In the normalized frame, the vertex is $s_0=(0,0,1)$.
For the first family, put
\[
 P_1=(\rho,-\rho,1),\quad P_2=(\rho,-2\rho,1),\quad
 P_3=(\rho,\rho,1+\rho/\sqrt\varepsilon).
\]
These are the vertices of $T_r$ when $r=\rho\sqrt\varepsilon/z_0$.
For \cref{dp:tangent-family}, put
\[
 P_1=(\rho,-\rho,1-2\rho(1-\eta)),\quad
 P_2=(\rho,-k\rho,1-2\sqrt k\rho(1-\eta)),\quad
 P_3=(\rho,\rho,1+\rho L),
\]
with $r=\rho/\zK$. All test points used below have nonnegative $q$.
Let $m_{01},m_{02},m_{12}$ be the three midpoints of $s_0,P_1,P_2$.
If a completed orbit set contains $T_r$, it must contain these six test
points and its generator must intersect the line
$\{(\rho,\rho,h):h\in\R\}$. These necessary conditions involve neither
$\varepsilon$ nor $L$. They apply even when the generating orbit set does
not contain $s_0$.

For a rational vector $v$, a point exclusion consists of
\begin{equation}\label{dp:point-exclusion}
 v^TGM(s)v<0,\qquad v^TGM(s-q(s)e_w)v<0.
\end{equation}
The same quadratic form is then negative throughout the allowable lowering
interval by affinity, so \cref{dp:completion-membership} excludes $s$.
For the vertical line, its PSD matrix has second diagonal entry
$\rho g_{21}+g_{22}$. Its determinant is
\[
 \det G(h-\rho^2)-\frac{(c-g_{21}h)^2}{4},\qquad
 c=\rho(g_{11}-g_{22})+g_{12}.
\]
If $g_{21}\ne0$, completing the square gives its maximum over $h$ as
\begin{equation}\label{dp:psi}
 \frac{\det G}{g_{21}^2}\Psi(G),\qquad
 \Psi(G)=\rho g_{21}(g_{11}-g_{22})+g_{11}g_{22}-\rho^2g_{21}^2.
\end{equation}
If $g_{21}=0$, $\Psi(G)=\det G>0$. Thus either
$\rho g_{21}+g_{22}<0$ or $\Psi(G)<0$ excludes the line whenever
$\det G>0$.

Normalize the largest absolute matrix entry to one by positive scaling.
This covers every nonzero $G$ by the eight facets
$g_{ij}=\pm1$, with the other three entries in $[-1,1]$.
Each facet is recursively bisected in a free coordinate with longest
interval, resolving ties by row-major order. A binary path records the
successive lower or upper half intervals. Each leaf has one of the following
exact exclusions:
\begin{enumerate}
 \item $\det G<0$ at every vertex of the box;
 \item a test point and rational vector satisfying both linear inequalities
       in \cref{dp:point-exclusion} throughout the box;
 \item an exact interval upper bound on $\Psi(G)$ that is negative;
 \item the upper bound of $\rho g_{21}+g_{22}$ that is negative.
\end{enumerate}
The first test is valid by multilinearity of the determinant in the free
entries. The second and fourth are linear box bounds. For the third, write
the intervals of $g_{11},g_{21},g_{22}$ as $I_a,I_c,I_d$. Let
$\max(I I')$ and $\min(I I')$ denote the maximum and minimum of the
four endpoint products. An upper bound is
\[
 \rho\max(I_cI_a)-\rho\min(I_cI_d)+\max(I_aI_d)
 -\rho^2\min_{c\in I_c}c^2,
\]
where the last minimum is zero if $0\in I_c$ and is the smaller squared
endpoint otherwise. All arithmetic is rational.

The complete leaf lists supplied with the evidence companion have the counts
in the table. In each facet their paths are prefix-free and satisfy
$\sum_{\text{leaves }p}2^{-|p|}=1$. A finite prefix-free family gives
disjoint box interiors; the exact sum gives coverage up to boundary, and
the closed boxes cover the boundary too. All exclusions are strict and
hold on those closed boxes. Thus every admissible normalized matrix is excluded from
at least one necessary condition. No completed set contains $T_r$ at the
listed threshold; monotonicity of $T_r$ then gives the stated supremum
upper bound.

\begin{center}
\begin{tabular}{@{}lrrrrr@{}}\toprule
 Family and $\rho$ & Leaves & Point & $\Psi$ & Det. & Diag.\\\midrule
 First family, $137$ & 28699 & 8664 & 19944 & 54 & 37\\
 $\eta=1/1000,k=4$, $197/200$ & 11132 & 10593 & 451 & 74 & 14\\
 $\eta=1/100,k=4$, $5/4$ & 1882 & 1690 & 150 & 27 & 15\\
 $\eta=1/1000,k=9/4$, $21/20$ & 10705 & 10300 & 267 & 76 & 62\\
 $\eta=1/1000,k=49/25$, $11/10$ & 11943 & 11499 & 245 & 93 & 106\\\bottomrule
\end{tabular}
\end{center}
This is a computer-assisted proof with exact certificates. The numerical
search used to discover the leaves is not part of the validity argument.
The companion includes the rational records and a verifier; an independent
verification also checked exact facet volumes, nonoverlap, and the exact
maximum of $\Psi$ on each relevant box. The hypotheses and exclusions above
are the complete mathematical reduction.

\subsection{Printed rational lower witnesses}\label{dp:lower-witnesses}

The lower bounds use the following matrices $G$, read as exact terminating
decimals, and rational heights $h_0$. Set $s_0=(0,0,1)$ and use the points
$P_1,P_2$ defined in \cref{dp:box-proof} at the stated $\rho$.

\begin{center}\small
\begin{tabular}{@{}lcl@{}}\toprule
 Family & $\rho$ & $G$ and $h_0$\\\midrule
 First family & $1367/10$ &
 $\begin{pmatrix}1&-1.678587\\0.007315&0.698332\end{pmatrix}$,
 $h_0=5417132036/169459$\\[6pt]
 $\eta=1/1000,k=4$ & $2461/2500$ &
 $\begin{pmatrix}1.00000003&-2.11841757\\1.01581925&0.30393078\end{pmatrix}$,
 $h_0=1240951/370546$\\[6pt]
 $\eta=1/100,k=4$ & $6117/5000$ &
 $\begin{pmatrix}1&-2.08684938\\0.81738698&0.4028837\end{pmatrix}$,
 $h_0=4011447/862162$\\[6pt]
 $\eta=1/1000,k=9/4$ & $5207/5000$ &
 $\begin{pmatrix}1&-2.43916573\\0.96020732&0.5468295\end{pmatrix}$,
 $h_0=4025363/954377$\\[6pt]
 $\eta=1/1000,k=49/25$ & $2693/2500$ &
 $\begin{pmatrix}1&-2.51367434\\0.9282616&0.62838339\end{pmatrix}$,
 $h_0=412614/89749$\\\bottomrule
\end{tabular}
\end{center}
The complete verification is four $2\times2$ rational matrix checks per row:
\begin{equation}\label{dp:witness-checks}
 \sym G\succ0,\qquad
 \sym(GM(P_1))\succeq0,\qquad
 \sym(GM(P_2))\succeq0,\qquad
 \sym(GM(\rho,\rho,h_0))\succeq0.
\end{equation}
For reproducibility without a matrix package, if
$G=\left(\begin{smallmatrix}a&b\\c&d\end{smallmatrix}\right)$ and
$s=(x,y,w)$, the three independent entries to test are
\[
 A=aw+by,\quad B=(ax+b+cw+dy)/2,\quad E=cx+d.
\]
PSD means $A,E\ge0$ and $AE-B^2\ge0$; PD means $A>0$ and
$AE-B^2>0$. Substitution of each printed row gives these signs exactly.
The finite exact arithmetic records in the companion contain the results.
Positive definiteness at $s_0$ also implies $\det G>0$.

Upward completion contains $(\rho,\rho,h)$ for every $h\ge h_0$.
In the first family, $1+\rho/\sqrt\varepsilon\ge h_0$ precisely when
$\varepsilon\le(\rho/(h_0-1))^2$, which is the exact $\varepsilon_1$
in \cref{dp:certified-sharpness}. In the four tangent families, the
individual thresholds $L\ge(h_0-1)/\rho$ are
\[
 \frac{1088006250}{455956853},\quad
 \frac{7873212500}{2636922477},\quad
 \frac{15354930000}{4969441039},\quad
 \frac{807162500}{241694057},
\]
all below $17/5$. The completed set then contains every vertex of $T_r$
and hence $T_r$. Its cut ratio is at least $r$, yielding every listed
lower bound. No optimization or assertion of optimality is needed for these
witnesses.

\subsection{A dual certificate for an orbit upper bound}\label{dp:dual-depth}

A second exact certificate illustrates the dual PSD obstruction and retains
an upper bound independent of the completed-set calculation. On the first
family in \cref{dp:vanishing}, it gives
\[
 \zA/\zK\le160\sqrt\varepsilon/z_0\quad
 (0<\varepsilon\le1/9765625).
\]
The completed-family bound in \cref{dp:certified-sharpness} is stronger.

Set $\rho=160$, $u=\sqrt\varepsilon/160$, and
\[
 M_1=\begin{pmatrix}1&160\\-160&1\end{pmatrix},\quad
 M_2=\begin{pmatrix}1&160\\-320&1\end{pmatrix},\quad
 M_3=\begin{pmatrix}1+1/u&160\\160&1\end{pmatrix}.
\]
Take
\[
 Y_1=\begin{pmatrix}979/10^6&-5679/500000\\-5679/500000&26393/200000\end{pmatrix},\quad
 Y_2=\begin{pmatrix}131/250000&859/100000\\859/100000&35253/250000\end{pmatrix},
\]
and
\[
 \begin{gathered}
 n=1/10,\quad \vartheta=140667062516/1953125,\quad
 \sigma_\infty=-187502/3125,\\
 \sigma(u)=\sigma_\infty+160\vartheta u^2,\quad
 Y_3(u)=\begin{pmatrix}\vartheta u^2&\sigma(u)u\\\sigma(u)u&n\end{pmatrix}.
 \end{gathered}
\]
Put $Y_0(u)=-(M_1Y_1+M_2Y_2+M_3Y_3(u))$. The sum is symmetric
identically. The fixed matrices $Y_1,Y_2$ are PD by their positive first
entries and determinants. The remaining signs can be checked without root
counting. For a polynomial $p(u)=\sum_k p_ku^k$ and
$0\le u\le u_*=1/500000$, use
\[
 p(u)\ge p_0+\sum_{k\ge1:\ p_k<0}p_ku_*^k.
\]
Direct expansion of the printed matrices gives the following exact lower
bounds on this right-hand expression:
\begin{center}\small
\begin{tabular}{@{}ll@{}}\toprule
 Polynomial & Lower bound\\\midrule
 $\vartheta n-\sigma(u)^2$ &
 $2096684649823859428966288734234359/582076609134674072265625000000>3600$\\
 $Y_0(u)_{11}$ &
 $24149614711567794247/76293945312500000000>3/10$\\
 $\det Y_0(u)$ &
 $176579577419584483919878545141452997/2328306436538696289062500000000000000>3/40$\\\bottomrule
\end{tabular}
\end{center}
Thus $Y_3(u),Y_0(u)\succ0$ for $0<u\le u_*$.
At $u=0$ the limiting matrix is
\[
 Y_0(0)=\begin{pmatrix}441377/10^6&20443/62500\\20443/62500&558543/10^6\end{pmatrix}.
\]
An admissible $G$ has $\sym G\succ0$ at the normalized vertex.
If it contained all three endpoint matrices, the trace identity
\[
 0=\tr\left(G\left(Y_0+\sum_{i=1}^3M_iY_i\right)\right)
   =\ip{\sym G}{Y_0}+\sum_{i=1}^3\ip{\sym(GM_i)}{Y_i}
\]
would have a strictly positive first term and nonnegative remaining terms.
This is impossible. The claimed upper bound follows from the printed
rational certificate and these elementary polynomial sign bounds.

\subsection{Derivation and sharpness of the fixed rule}\label{dp:fixed-rule-derivation}

In Sylvester coordinates,
\[
 M=x_1I+x_2J+y_1\diag(1,-1)+y_2
       \begin{pmatrix}0&1\\1&0\end{pmatrix},\quad
 (x_1,x_2,y_1,y_2)=\tfrac12(w+1,x-y,w-1,x+y).
\]
Let $\sin\theta=(\bar x-\bar y)/N_0$ and
$\cos\theta=(\bar w+1)/N_0$. Multiplication by
\[
 R_\theta=\begin{pmatrix}\cos\theta&-\sin\theta\\
                         \sin\theta&\cos\theta\end{pmatrix}
\]
rotates the positive coordinates so that
$\sym(R_\theta M)\succeq0$ is the chosen inequality
$\lambda^Tx\ge\norm y$. Its determinant is one and its skew affine form is
\[
 \ell=\cos\theta(x-y)-\sin\theta(w+1)
     =-V(s-\sbar)/N_0.
\]
\Cref{dp:orbit-equation} now gives \cref{dp:fixed-rule-set}.
Because the vertex has $q_0>0$ and $\ell(\sbar)=0$ with positive trace,
it is positive definite. Along a ray, the first determinant zero is the
uncompleted exit. The completed exit is instead determined by
\cref{dp:lowering}. For example, at $\sbar=(1,0,1)$ on $p=e_w$,
the uncompleted determinant equation is $20(1+s)=s^2$, giving
$s=10+\sqrt{120}$. The completion contains the entire upward ray.

On the first sharp family, $\bar x=\bar y=0$ makes the fixed set the
upward-closed cylinder $w\ge(x+y)^2/4$. Its ray steps are
$\infty,2\sqrt\varepsilon,z_0$. Therefore
$z_{\mathrm{fix}}=\min\{2\sqrt\varepsilon,z_0\}$, with
$2\sqrt\varepsilon\le z_0$ exactly when $\varepsilon\le4/9$.

\begin{proposition}[A pure rescaling obstruction]\label{dp:rescaling}
Let $L\ge1$, $\sbar=(0,0,1)$,
$p_1=(L,-1/L,0)$, $p_2=(L,1/L,0)$, $p_3=e_w$, and unit costs.
Then $\zK=1$, $D=1$, the discriminants are $-1,+1,+1$, and
\[
 \cond(P)=\begin{cases}\sqrt2L,&1\le L\le\sqrt2,\\L^2,&L\ge\sqrt2.\end{cases}
\]
The orbit cylinder with $t=1/L$ is exact, but
$z_{\mathrm{fix}}=2/(L+1/L)$.
\end{proposition}

\begin{proof}
Expansion gives $q=1+\lambda_3+\lambda_1^2-\lambda_2^2$, so objective
one is feasible only on ray two, with transversal contact.
The cylinder $q\ge(x/L-Ly)^2/4$ contains $T^*$ and has the vertex inside.
The fixed cylinder $w\ge(x+y)^2/4$ exits on ray two at
$2/(L+1/L)$, no later than on ray one. Its completion is the same cylinder.
The rows of $P$ are orthogonal with norms $\sqrt2L,\sqrt2/L,1$,
giving the stated condition number. The family is the image of $L=1$
under $(x,y,w)\mapsto(Lx,y/L,w)$.
\end{proof}

\subsubsection{Proof of the prescribed condition number and depth example}\label{dp:fixed-rule-sharp-proof}

For the construction in \cref{dp:fixed-rule-sharp}, expansion gives
\[
 q(\sbar+P\lambda)=v^2(1-\lambda_3)+uvD(\lambda_1+\lambda_2)
                          +v^2D^2\lambda_1\lambda_2.
\]
Thus $\zK=1$, uniquely on ray three. The restrictions on the first two
rays are $v^2+uvDs$ and on the third $v^2(1-s)$, proving the extreme
relative discriminants. Also $q_0=v^2$, $\widetilde X=\widetilde Y=vD$,
and the singular values are $vD,vD,v^2$. Since $v=\kappa D$,
their ratio is exactly $\kappa$.

Translate the contact point $t^*=(u,-u,-u^2)$ to zero by
\[
 (x',y',w')=(x-u,y+u,w+ux-uy-u^2).
\]
The other vertices become $(0,0,v^2)$,
$(vD,0,v^2+uvD)$, and $(0,-vD,v^2+uvD)$.
Use \cref{dp:contact-set} with $\alpha=1$ and
$\beta=D/(uD+v)$. At the vertex its determinant slack is
$v^2(4-\beta^2v^2)>0$, since $\beta v<1$. At the two other vertices
the squared affine form vanishes, and their slack is
$4(v^2+uvD)>0$. At the contact the slack is zero.
The trace branch $w'+\beta x'+1$ is positive at all four vertices.
This orbit set therefore contains $T^*$ with $\sbar$ inside, proving
attained exactness for both families.

For the fixed rule, $\bar x-\bar y=2u$ and $\bar w+1=0$, so
$C_{\mathrm{fix}}=\{4q\ge(w+1)^2,\ x\ge y\}$.
The first two rays stay inside forever. On ray three the endpoint solves
$4v^2(1-s)=v^4s^2$, giving $s=2/(u+1)$.
The vertical section at $(u,-u)$ is
$w\in[1-2u,1+2u]$. Its upward completion is
$w\ge1-2u$, so the downward ray has exactly the same step.
This proves the claimed completed and uncompleted values.

The two-variable fixed-rule example in \cref{fd:fixed-rule} does not settle
the analogous margin or conditioning question for Case~2. Its horizontal
ray restriction is $(x_0-s)^2+1$, with relative discriminant
$-4/(8x_0^2+4)\to0$. Its two scaled ray lengths are $\zK/\epsilon$ and
$\zK$, so their condition number is $1/\epsilon$. Both proposed safeguards
therefore fail along its vanishing-ratio sequence.

\subsection{Full details of the contact approximation}\label{dp:contact-proof}

Let $A_0=\sym(GM(\sbar))$ and $A_i=\sym(GM(v_i))$ be PSD.
For fixed $0<r<1$, put
$A_i(r)=(1-r)A_0+rA_i$. If $z\in\ker A_i(r)$, its quadratic form
is a positive combination of two nonnegative quadratic forms, so both
vanish. For a PSD matrix $A$, $z^TAz=0$ implies $Az=0$, by diagonalizing
$A$. Hence
\[
 \ker A_i(r)=\ker A_0\cap\ker A_i.
\]
Such a vector is annihilated by the symmetric part at every point of the
segment from $\sbar$ to $v_i$. On the shortened segment to $v_i(r)$,
the objective level is below $\zK$, so $q>0$ throughout. Write
$GM(s)=A(s)+k(s)J$. If the common kernel is nonzero, then
$GM(s)z=k(s)Jz$ for a nonzero $z$. The left side cannot vanish, since
$\det G>0$ and $\det M(s)=q(s)>0$. Thus $k(s)\ne0$ and its continuous
sign is constant. Comparing the identities at the two shortened endpoints
gives
\[
 M(v_i(r))z=\frac{k(v_i(r))}{k(\sbar)}M(\sbar)z.
\]
Consequently the symmetric perturbation
$H_i=\sym(M(\sbar)^{-1}M(v_i(r)))$ satisfies
\[
 z^TH_i z=\frac{k(v_i(r))}{k(\sbar)}\norm z^2>0
 \quad\text{on }\ker A_i(r)\setminus\{0\}.
\]

Here is the perturbation argument without a limiting eigenvalue assertion.
In an orthogonal decomposition into the positive eigenspace and kernel of
$A_i(r)$, its perturbation has blocks
\[
 \begin{pmatrix}A_++\epsilon H_{++}&\epsilon H_{+0}\\
                  \epsilon H_{0+}&\epsilon H_{00}\end{pmatrix}.
\]
The first block remains positive definite for small $\epsilon$.
Its Schur complement is
\[
 \epsilon H_{00}-\epsilon^2H_{0+}
 (A_++\epsilon H_{++})^{-1}H_{+0},
\]
which is positive definite for all sufficiently small positive $\epsilon$
because $H_{00}\succ0$ and the inverse is bounded. If there is no kernel,
ordinary continuity suffices. If the whole matrix is zero, the perturbation
itself is positive definite by the displayed kernel inequality.
There are finitely many vertices, so the minimum of their positive admissible
epsilon thresholds is positive. Also $\det(G+\epsilon M(\sbar)^{-1})>0$
for small $\epsilon$. At $\sbar$ the symmetric matrix is $A_0+\epsilon I$.
These facts prove \cref{dp:contact-approx} completely.

For completeness, the singular-limit necessity in
\cref{dp:contact-intervals} uses uniqueness rather than dimension.
Normalize a sequence with cut ratios tending to one to Frobenius norm one.
Closed PSD inequalities pass to a nonzero limit $G$ on every vertex of
$T^*$, and $\det G\ge0$. If $\det G=0$, membership at zero gives
$G=\left(\begin{smallmatrix}a&0\\c&d\end{smallmatrix}\right)$ with
$d\ge0$, $ad=0$.
If $d>0$, the PSD matrix has zero first diagonal entry, so its
off-diagonal entry $(ch+dy)/2$ vanishes. Hence $y=-(c/d)h$.
Set $\beta=c/d$ and take
\[
 0<\alpha<\min_{v\ne0,\ x_v\ne0}\frac{4q(v)}{x_v^2},
\]
with an empty minimum interpreted as infinity. Its determinant slack is
strictly positive at every noncontact vertex.
If $d=0$, PSD membership at $\sbar$, whose height is positive, forces
$a\ge0$. If $a=0$, the zero first diagonal entry would force
$ch_{\sbar}=0$, and then $G=0$, impossible. Thus $a>0$.
The determinant of its symmetric matrix is
$-(ax-ch)^2/4$, so $x=(c/a)h$ at every vertex. Choose
\[
 \alpha>\max_{v\ne0}\frac{y_v^2}{4q(v)},\qquad\beta=\alpha c/a.
\]
Again the determinant slack is strict. In both cases $q(v)>0$ at every
noncontact vertex follows from uniqueness, and finite cardinality gives a
single parameter. Their first diagonal entries $\alpha h_v$ are positive,
so the strict determinant slack implies PD membership. The contact itself
has PSD matrix $e_2e_2^T$. Thus a singular limit actually leads to an
attaining set; it cannot obstruct closed-interval feasibility.

\subsection{The explicit threshold counterexample}\label{dp:cylinder-sharp-proof}

At $d=0$, write $a,b,c$ for the barycentric coordinates on the face opposite
the contact, in the order $\sbar,v_2,v_3$. Substituting $b=1-a-c$ gives
\[
 q_0=2a^2+\tfrac{11}{2}ac-\tfrac54a+\tfrac54c^2-\tfrac12c+\tfrac14.
\]
Its Hessian has determinant $-81/4<0$, so an interior stationary point
cannot minimize it. On $c=0$, the polynomial
$2a^2-5a/4+1/4$ has minimum $7/128$ at $a=5/16$.
On $a=0$ the minimum of $5c^2/4-c/2+1/4$ is $1/5$.
On $b=0$ the expression is $1+9a(1-a)/4$, with minimum one.
This proves the face bound used in \cref{dp:cylinder-sharp}.
The determinant of the ray matrix is $-3(3d+4)/2$, so these instances also
have full-dimensional optimal simplices.

It remains to exclude closed interval intersections. Put $\alpha=t^2$.
The three vertex $q$ values are $1-d,1/4-d,1-d$. Let $L_0,L_3$ be the
left endpoints of the intervals for $\sbar,v_3$, and $R_2$ the right
endpoint for $v_2$:
\begin{align*}
 L_0(t)&=\frac{3t^2/2+1/2-2t\sqrt{1-d}}{1/4-d},\\
 R_2(t)&=\frac{t^2/2-3/2+2t\sqrt{1/4-d}}{1-d},\\
 L_3(t)&=\frac{3t^2-1-2t\sqrt{1-d}}{4-d}.
\end{align*}
For $0<d\le1/8$,
$1-\sqrt{1-d}\ge d/2$ and
$\sqrt{1/4-d}\le(1-2d)/2$.
If $t\ge1$, the numerator of $L_0$ is at least
$(3t-1)(t-1)/2+td\ge0$. Hence
\[
 L_0\ge2(3t-1)(t-1)+4td,\qquad
 R_2\le\tfrac47(t+3)(t-1),
\]
where the second bound discards a negative term and uses $1/(1-d)\le8/7$.
Thus
\[
 L_0-R_2\ge(t-1)(38t-26)/7+4dt>0.
\]
If $0<t\le1$, split the numerator of $L_3$ as
$(3t+1)(t-1)+2t(1-\sqrt{1-d})$.
The first summand is nonpositive and the second nonnegative, so
\[
 L_3\ge\tfrac8{31}(3t+1)(t-1)+td/4.
\]
The numerator of $R_2$ is at most
$(t+3)(t-1)/2-2td\le0$, giving
$R_2\le(t+3)(t-1)/2-2td$.
Therefore
\[
 L_3-R_2\ge(1-t)(77-17t)/62+(9/4)dt>0.
\]
In each case two closed intervals are disjoint. This proves the strict
orbit gap throughout $0<d<7/128$. Taking $d$ arbitrarily small proves
the sharp cylinder threshold analytically; finite checks at three values
are not needed for that conclusion.

\subsection{Saved numerical calibration and its exact correction}\label{dp:calibration}

The exact bounds above distinguish certification from numerical set
optimization. On the fixed-ray family, normalized numerical optimization of
the uncompleted family suggests a leading coefficient about $56.7$ in
$\sqrt\varepsilon/z_0$. The completed-family search found about $136.707$.
Only the rational interval $[136.7,137]$ is used as a proved completed-family
coefficient in the stated small-$\varepsilon$ range. Similarly, the
construction that lowers the third vertex proves a comparison with the best
finite-height uncompleted set, but does not prove that all optimal completed
sets arise this way.
Each printed lower witness has its generating orbit set containing the
vertex in its interior. On the first family, restricting a search to such
generators therefore loses at most $137/(1367/10)-1<0.22\%$ in the stated
small-$\varepsilon$ range. This estimate is specific to that family.

A retained near-boundary corner also illustrates the need to normalize
before treating numerical infeasibility as evidence. Its initial
original-coordinate optimization underestimated the orbit ratio. Rational
primal and dual certificates instead give
\[
 \frac{763}{25000}=0.03052\ \le\ \frac{\zA}{\zK}\ \le\
 \frac{191}{6250}=0.03056.
\]
The retained verification records fix the input as the exact binary values
of the stored floating-point data. They report a rational primal check by
PSD membership with the vertex PD. The reported dual check uses
the trace identity in \cref{dp:dual-depth}, with PSD endpoint matrices and
a PD vertex matrix. The objective normalization uses the exact corner
bound of that specified input. The record reports an exact bracket, rather
than a numerical solver bound. Its full rational primal and dual matrices
were not archived; the companion retains the checking source and its
successful exact-arithmetic receipt. This historical correction has that
reproducibility limitation and is not used in any theorem above.
